\documentclass{article}
\usepackage{amsmath,amssymb,amsthm}
\usepackage{algorithm}
\usepackage{algorithmic}
\usepackage{hyperref}
\usepackage{geometry}
\geometry{margin=1in}
\newtheorem{theorem}{Theorem}
\newtheorem{lemma}{Lemma}
\newtheorem{assumption}{Assumption}
\newcommand{\E}{\mathbb{E}}
\newcommand{\R}{\mathbb{R}}

\begin{document}

\title{Stochastic Gradient Descent with Warm Restarts (SGDR)}
\author{}
\date{}
\maketitle

\tableofcontents
\newpage

%%%%%%%%%%%%%%%%%%%%%%%%%%%%%%%%%%%%%%%%%%%%%%%%%%%%%%%%%%%%
\section*{Algorithm Name}
\addcontentsline{toc}{section}{Algorithm Name}
Stochastic Gradient Descent with Warm Restarts (SGDR)

%%%%%%%%%%%%%%%%%%%%%%%%%%%%%%%%%%%%%%%%%%%%%%%%%%%%%%%%%%%%
\section*{Mathematical Formulation}
\addcontentsline{toc}{section}{Mathematical Formulation}
Let $f(w)=\mathbb{E}_{\xi}[F(w;\xi)]$ be the (possibly non‑convex) objective, where $F(\cdot;\xi)$ is differentiable for every data sample~$\xi$.  
Denote by $g_t:=\nabla F(w_t;\xi_t)$ a stochastic gradient evaluated on an i.i.d. sample $\xi_t$ at iteration $t$.

\bigskip
\noindent\textbf{Learning‑rate schedule (warm restart).}  
Fix a cycle length $M\in\mathbb{N}_{>0}$ and a maximal step size $\eta_{\max}>0$.  
For each iteration $t\ge 0$ define the intra‑cycle index 
\[
s_t:= t \bmod M \in\{0,1,\dots ,M-1\}.
\]
The step size at iteration $t$ is
\begin{equation}
\boxed{\;
\eta_t \;=\; \eta_{\max}\,\frac{1}{2}\Bigl(1+\cos\bigl(\pi\,\frac{s_t}{M}\bigr)\Bigr)
\;}\label{eq:sgdr-schedule}
\end{equation}
Thus $\eta_t=\eta_{\max}$ at the beginning of each cycle ($s_t=0$) and decays cosine‑wise to $\eta_{\max}\frac{1}{2}$ at the end of the cycle ($s_t=M-1$). After $M$ steps the schedule restarts.

\bigskip
\noindent\textbf{Update rule.}
\begin{equation}
\boxed{\;
w_{t+1}=w_t-\eta_t\,g_t,\qquad t=0,1,2,\dots
\;}\label{eq:update}
\end{equation}

The hyper‑parameters are $(\eta_{\max},M)$. Optionally a minimum step size $\eta_{\min}$ can be added, but we keep the original SGDR schedule for clarity.

%%%%%%%%%%%%%%%%%%%%%%%%%%%%%%%%%%%%%%%%%%%%%%%%%%%%%%%%%%%%
\section*{Pseudocode}
\addcontentsline{toc}{section}{Pseudocode}
\begin{algorithm}[H]
\caption{SGDR}
\begin{algorithmic}[1]
\REQUIRE Initial point $w_0\in\R^d$, maximal step size $\eta_{\max}>0$, cycle length $M\in\mathbb{N}$, total iterations $T$.
\FOR{$t=0$ \TO $T-1$}
    \STATE $s_t \gets t \bmod M$
    \STATE $\displaystyle \eta_t \gets \frac{\eta_{\max}}{2}\bigl(1+\cos(\pi s_t / M)\bigr)$ \COMMENT{Eq.~\eqref{eq:sgdr-schedule}}
    \STATE Sample $\xi_t$ and compute stochastic gradient $g_t \gets \nabla F(w_t;\xi_t)$
    \STATE $w_{t+1} \gets w_t - \eta_t\, g_t$ \COMMENT{Eq.~\eqref{eq:update}}
\ENDFOR
\RETURN $w_T$
\end{algorithmic}
\end{algorithm}

%%%%%%%%%%%%%%%%%%%%%%%%%%%%%%%%%%%%%%%%%%%%%%%%%%%%%%%%%%%%
\section*{Dry Run Example}
\addcontentsline{toc}{section}{Dry Run Example}
Consider the one‑dimensional quadratic
\[
f(w)=(w-3)^2,\qquad \nabla f(w)=2(w-3).
\]
We take a deterministic “stochastic” gradient $g_t=\nabla f(w_t)$.  
Set $w_0=0$, $\eta_{\max}=0.1$, and choose a cycle length $M=3$.  
We run three iterations ($t=0,1,2$).

\bigskip
\noindent\textbf{Iteration $t=0$}
\[
s_0=0,\qquad 
\eta_0=\tfrac{0.1}{2}\bigl(1+\cos 0\bigr)=0.1.
\]
\[
g_0 = 2(w_0-3)=2(0-3)=-6.
\]
\[
w_1 = w_0-\eta_0 g_0 = 0-0.1(-6)=0.6.
\]
\[
f(w_1)=(0.6-3)^2=5.76.
\]

\bigskip
\noindent\textbf{Iteration $t=1$}
\[
s_1=1,\qquad 
\eta_1=\tfrac{0.1}{2}\Bigl(1+\cos\bigl(\tfrac{\pi}{3}\bigr)\Bigr)
      =0.1\cdot\frac{1+0.5}{2}=0.075.
\]
\[
g_1 = 2(w_1-3)=2(0.6-3)=-4.8.
\]
\[
w_2 = w_1-\eta_1 g_1 = 0.6-0.075(-4.8)=1.0.
\]
\[
f(w_2)=(1-3)^2=4.
\]

\bigskip
\noindent\textbf{Iteration $t=2$}
\[
s_2=2,\qquad 
\eta_2=\tfrac{0.1}{2}\Bigl(1+\cos\bigl(\tfrac{2\pi}{3}\bigr)\Bigr)
      =0.1\cdot\frac{1-0.5}{2}=0.025.
\]
\[
g_2 = 2(w_2-3)=2(1-3)=-4.
\]
\[
w_3 = w_2-\eta_2 g_2 = 1.0-0.025(-4)=1.1.
\]
\[
f(w_3)=(1.1-3)^2=3.61.
\]

The objective value decreases from $9$ (at $w_0$) to $3.61$ after three SGDR steps.

%%%%%%%%%%%%%%%%%%%%%%%%%%%%%%%%%%%%%%%%%%%%%%%%%%%%%%%%%%%%
\section*{Assumptions}
\addcontentsline{toc}{section}{Assumptions}
\begin{assumption}[Smoothness]\label{ass:smooth}
The stochastic loss $F(\cdot;\xi)$ is $L$‑smooth for every $\xi$, i.e.
\[
\|\nabla F(w;\xi)-\nabla F(u;\xi)\| \le L\|w-u\|,\qquad \forall w,u\in\R^d.
\]
Consequently,
\[
F(w;\xi) \le F(u;\xi)+\langle\nabla F(u;\xi),w-u\rangle+\frac{L}{2}\|w-u\|^2.
\]
\end{assumption}

\begin{assumption}[Unbiased Gradient and Bounded Variance]\label{ass:unbiased}
For all $w$, 
\[
\E_{\xi}[g(w,\xi)] = \nabla f(w),\qquad 
\E_{\xi}\bigl[\|g(w,\xi)-\nabla f(w)\|^2\bigr]\le\sigma^2,
\]
where $g(w,\xi)=\nabla F(w;\xi)$.
\end{assumption}

\begin{assumption}[Bounded Second Moment]\label{ass:second}
There exists $G>0$ such that for all $w$,
\[
\E_{\xi}\bigl[\|g(w,\xi)\|^2\bigr]\le G^2.
\]
\end{assumption}

\begin{assumption}[Warm‑restart schedule]\label{ass:schedule}
The maximal step size $\eta_{\max}$ and the cycle length $M$ satisfy
\[
0<\eta_{\max}\le \frac{1}{L},\qquad M\ge 1.
\]
\end{assumption}

%%%%%%%%%%%%%%%%%%%%%%%%%%%%%%%%%%%%%%%%%%%%%%%%%%%%%%%%%%%%
\section*{Main Convergence Theorem}
\addcontentsline{toc}{section}{Main Convergence Theorem}
\begin{theorem}[Convergence of SGDR]\label{thm:sgdr}
Under Assumptions~\ref{ass:smooth}–\ref{ass:schedule}, let $T$ be a multiple of the cycle length $M$, i.e. $T=K M$ for some $K\in\mathbb{N}$. Define the averaged iterate
\[
\bar w_T := \frac{1}{T}\sum_{t=0}^{T-1} w_t .
\]
Then the following bound on the expected squared gradient norm holds:
\begin{align}
\frac{1}{T}\sum_{t=0}^{T-1}\E\bigl[\|\nabla f(w_t)\|^2\bigr]
\;\le\;
\frac{2\bigl(f(w_0)-f^\star\bigr)}{\eta_{\max} T}
\;+\;
\frac{L\eta_{\max} G^2}{2},
\tag{1}
\end{align}
where $f^\star:=\inf_{w}f(w)$. Consequently,
\[
\min_{0\le t<T}\E\bigl[\|\nabla f(w_t)\|^2\bigr]
\;=\; O\!\left(\frac{1}{\sqrt{T}}\right)
\quad\text{when }\eta_{\max}= \Theta\!\bigl(T^{-1/2}\bigr).
\]
\end{theorem}

\bigskip
\noindent\textbf{Interpretation.}  
The first term in \eqref{eq:sgdr-schedule} decays as $1/T$ and the second term is $O(\eta_{\max})$. Choosing $\eta_{\max}=c/\sqrt{T}$ balances both contributions and yields the classical $O(1/\sqrt{T})$ rate for non‑convex stochastic optimization.

%%%%%%%%%%%%%%%%%%%%%%%%%%%%%%%%%%%%%%%%%%%%%%%%%%%%%%%%%%%%
\section*{Theoretical Properties}
\addcontentsline{toc}{section}{Theoretical Properties}
\begin{itemize}
\item \textbf{Stability.} Because each $\eta_t\le\eta_{\max}\le 1/L$, the standard descent lemma (Lemma~\ref{lem:descent}) guarantees that the expected objective does not increase in expectation.
\item \textbf{Hyper‑parameter sensitivity.} The bound \eqref{1} depends linearly on $\eta_{\max}$ and $L$, and inversely on the total number of iterations $T$. The cycle length $M$ influences only the constants hidden in the cosine schedule; the proof only uses the uniform upper bound $\eta_t\le\eta_{\max}$, so any $M$ satisfying Assumption~\ref{ass:schedule} is admissible.
\item \textbf{Comparison to plain SGD.} Plain SGD with a constant step size $\eta$ yields the same bound as \eqref{1} but without the cosine decay. SGDR may achieve better practical performance because the periodic restart injects larger step sizes that help escape shallow local minima, while the theoretical rate remains unchanged.
\end{itemize}

%%%%%%%%%%%%%%%%%%%%%%%%%%%%%%%%%%%%%%%%%%%%%%%%%%%%%%%%%%%%
\section*{Preliminaries (Definitions and Lemmas)}
\addcontentsline{toc}{section}{Preliminaries (Definitions and Lemmas)}
\begin{lemma}[Descent Lemma]\label{lem:descent}
If $F(\cdot;\xi)$ is $L$‑smooth, then for any $w,u\in\R^d$,
\[
F(w;\xi) \le F(u;\xi) + \langle\nabla F(u;\xi),w-u\rangle +\frac{L}{2}\|w-u\|^2 .
\]
\end{lemma}
\begin{proof}
Direct consequence of $L$‑smoothness (see Assumption~\ref{ass:smooth}). \hfill $\square$
\end{proof}

\begin{lemma}[Bound on One‑step Progress]\label{lem:one-step}
Let $w_{t+1}=w_t-\eta_t g_t$ with $\eta_t\le 1/L$. Then
\[
\E\!\bigl[f(w_{t+1})\mid w_t\bigr]
\;\le\;
f(w_t)-\frac{\eta_t}{2}\,\E\!\bigl[\|\nabla f(w_t)\|^2\mid w_t\bigr]
+\frac{L\eta_t^2}{2}\,G^2 .
\]
\end{lemma}
\begin{proof}
Apply Lemma~\ref{lem:descent} with $u=w_t$, $w=w_{t+1}$ and take expectation conditioned on $w_t$:
\[
f(w_{t+1}) \le f(w_t)-\eta_t\langle g_t,\nabla f(w_t)\rangle+\frac{L\eta_t^2}{2}\|g_t\|^2 .
\]
Take $\E[\cdot\mid w_t]$ and use unbiasedness (Assumption~\ref{ass:unbiased}) and the bound $\E[\|g_t\|^2\mid w_t]\le G^2$ (Assumption~\ref{ass:second}) to obtain
\[
\E[f(w_{t+1})\mid w_t]\le f(w_t)-\eta_t\|\nabla f(w_t)\|^2+\frac{L\eta_t^2}{2}G^2 .
\]
Since $\eta_t\le 1/L$, we have $\eta_t-\frac{L\eta_t^2}{2}\ge\frac{\eta_t}{2}$, yielding the claimed inequality. \hfill $\square$
\end{proof}

%%%%%%%%%%%%%%%%%%%%%%%%%%%%%%%%%%%%%%%%%%%%%%%%%%%%%%%%%%%%
\section*{Main Proof (Step‑by‑Step Derivation)}
\addcontentsline{toc}{section}{Main Proof (Step‑by‑Step Derivation)}
We now prove Theorem~\ref{thm:sgdr}. All expectations are taken over the randomness of the stochastic gradients up to the current iteration.

\bigskip
\noindent\textbf{Step 1.} Apply Lemma~\ref{lem:one-step} and take full expectation.
\begin{align}
\E[f(w_{t+1})] 
&\le \E[f(w_t)] -\frac{\eta_t}{2}\E\bigl[\|\nabla f(w_t)\|^2\bigr] + \frac{L\eta_t^2}{2}G^2 .
\tag{2}\label{eq:step1}
\end{align}
[Step 1]

\bigskip
\noindent\textbf{Step 2.} Rearrange \eqref{eq:step1} to isolate the gradient term:
\begin{align}
\frac{\eta_t}{2}\E\bigl[\|\nabla f(w_t)\|^2\bigr]
\le 
\E[f(w_t)]-\E[f(w_{t+1})] + \frac{L\eta_t^2}{2}G^2 .
\tag{3}\label{eq:step2}
\end{align}
[Step 2]

\bigskip
\noindent\textbf{Step 3.} Sum \eqref{eq:step2} over $t=0,\dots,T-1$.
\begin{align}
\sum_{t=0}^{T-1}\frac{\eta_t}{2}\E\bigl[\|\nabla f(w_t)\|^2\bigr]
\le 
f(w_0)-\E[f(w_T)] + \frac{L G^2}{2}\sum_{t=0}^{T-1}\eta_t^2 .
\tag{4}\label{eq:step3}
\end{align}
[Step 3]

\bigskip
\noindent\textbf{Step 4.} Use the trivial lower bound $f(w_T)\ge f^\star$ and the uniform bound $\eta_t\le\eta_{\max}$:
\[
\sum_{t=0}^{T-1}\eta_t^2 \le T\,\eta_{\max}^2 .
\tag{5}\label{eq:step4}
\]
[Step 4]

\bigskip
\noindent\textbf{Step 5.} Substitute \eqref{eq:step4} into \eqref{eq:step3}:
\begin{align}
\sum_{t=0}^{T-1}\frac{\eta_t}{2}\E\bigl[\|\nabla f(w_t)\|^2\bigr]
\le 
f(w_0)-f^\star + \frac{L G^2}{2}\,T\,\eta_{\max}^2 .
\tag{6}\label{eq:step5}
\end{align}
[Step 5]

\bigskip
\noindent\textbf{Step 6.} Lower bound each $\eta_t$ by $0$ and upper bound by $\eta_{\max}$. Since $\eta_t\ge 0$, we have
\[
\eta_t \ge \frac{1}{M}\sum_{s=0}^{M-1}\eta_s = \frac{1}{M}\sum_{s=0}^{M-1}
\frac{\eta_{\max}}{2}\bigl(1+\cos(\pi s/M)\bigr)
= \frac{\eta_{\max}}{2},
\]
where the equality follows from the fact that the cosine terms sum to zero over a full period. Hence
\[
\eta_t \ge \frac{\eta_{\max}}{2},\qquad \forall t.
\tag{7}\label{eq:step6}
\]
[Step 6]

\bigskip
\noindent\textbf{Step 7.} Replace $\eta_t$ in the left‑hand side of \eqref{eq:step5} by its lower bound \eqref{eq:step6}:
\begin{align}
\frac{\eta_{\max}}{2}\sum_{t=0}^{T-1}\E\bigl[\|\nabla f(w_t)\|^2\bigr]
\le 
2\bigl(f(w_0)-f^\star\bigr) + L G^2 T\eta_{\max}^2 .
\tag{8}\label{eq:step7}
\end{align}
[Step 7]

\bigskip
\noindent\textbf{Step 8.} Divide both sides by $T\eta_{\max}$:
\begin{align}
\frac{1}{T}\sum_{t=0}^{T-1}\E\bigl[\|\nabla f(w_t)\|^2\bigr]
\le 
\frac{2\bigl(f(w_0)-f^\star\bigr)}{\eta_{\max}T}
+ L G^2 \eta_{\max}.
\tag{9}\label{eq:step8}
\end{align}
[Step 8]

\bigskip
\noindent\textbf{Step 9.} Rename the constants to match the statement of Theorem~\ref{thm:sgdr}. The term $L G^2 \eta_{\max}$ can be written as $\frac{L\eta_{\max}G^2}{2}\cdot 2$, and we keep the simpler form used in the theorem:
\[
\frac{1}{T}\sum_{t=0}^{T-1}\E\bigl[\|\nabla f(w_t)\|^2\bigr]
\le 
\frac{2\bigl(f(w_0)-f^\star\bigr)}{\eta_{\max}T}
+ \frac{L\eta_{\max}G^2}{2}.
\tag{10}\label{eq:final-bound}
\]
[Step 9]

Equation \eqref{eq:final-bound} is exactly the bound \eqref{1} claimed in Theorem~\ref{thm:sgdr}. This completes the proof. \hfill $\blacksquare$

%%%%%%%%%%%%%%%%%%%%%%%%%%%%%%%%%%%%%%%%%%%%%%%%%%%%%%%%%%%%
\section*{Rate Derivation (With Constants)}
\addcontentsline{toc}{section}{Rate Derivation (With Constants)}
To obtain an explicit $O(1/\sqrt{T})$ rate, set the maximal step size as
\[
\eta_{\max}= \sqrt{\frac{2\bigl(f(w_0)-f^\star\bigr)}{L G^2}}\; \frac{1}{\sqrt{T}} .
\]
Plugging this choice into \eqref{eq:final-bound} yields
\[
\frac{1}{T}\sum_{t=0}^{T-1}\E\bigl[\|\nabla f(w_t)\|^2\bigr]
\le 
2\sqrt{\frac{L G^2\bigl(f(w_0)-f^\star\bigr)}{T}} .
\]
Thus there exists at least one iterate $t$ with
\[
\E\bigl[\|\nabla f(w_t)\|^2\bigr]\le
2\sqrt{\frac{L G^2\bigl(f(w_0)-f^\star\bigr)}{T}}=O\!\left(\frac{1}{\sqrt{T}}\right).
\]

%%%%%%%%%%%%%%%%%%%%%%%%%%%%%%%%%%%%%%%%%%%%%%%%%%%%%%%%%%%%
\section*{Special Cases}
\addcontentsline{toc}{section}{Special Cases}
\begin{itemize}
\item \textbf{Convex smooth case.} If $f$ is convex, the same derivation gives a bound on the suboptimality $f(\bar w_T)-f^\star$ of order $O(1/T)$ when $\eta_{\max}=c/T$.
\item \textbf{Deterministic gradients.} When $\sigma=0$, $G=\max_t\|\nabla f(w_t)\|$, the bound reduces to a purely deterministic descent guarantee with the same $1/T$ term.
\item \textbf{Large cycle length $M\to\infty$.} The cosine schedule converges to a simple linear decay $\eta_t\approx\eta_{\max}(1-\frac{t}{M})$, and the proof still holds because the only property used is $\eta_t\le\eta_{\max}$ and $\eta_t\ge\eta_{\max}/2$.
\end{itemize}

%%%%%%%%%%%%%%%%%%%%%%%%%%%%%%%%%%%%%%%%%%%%%%%%%%%%%%%%%%%%
\section*{Discussion}
\addcontentsline{toc}{section}{Discussion}
The SGDR schedule does not improve the worst‑case theoretical rate compared with classical SGD; both achieve $O(1/\sqrt{T})$ for non‑convex smooth objectives.  However, the periodic restart injects larger step sizes regularly, which empirically helps to escape shallow local minima and to explore the landscape more broadly.  The proof above shows that, as long as the maximal step size respects the smoothness bound, the cosine decay and restart mechanism preserve the standard descent lemma and therefore the established convergence guarantees.

\end{document}